\documentclass[11pt]{article}
\usepackage[T1]{fontenc}
\usepackage{amsmath,amssymb}
\usepackage{hyperref}

\title{Compatibility Findings: Classical Mechanics and FakeScience}
\author{Formalization notes}
\date{\today}

\begin{document}
\maketitle

\section{Findings}

The two supplied Lean files are not jointly importable as written. Both
declare the fully qualified name
\texttt{ClassicalMechanics.secondLawHolds}, but give it different types.
In \texttt{ClassicalMechanics.lean}, it states, for a system $M$ and a
particle $p$ at time $t$,
\[
  \operatorname{forceSum}(M.\mathrm{force}(p,t))
  = \operatorname{scale}(M.\mathrm{mass}(p),M.\mathrm{acceleration}(p,t)).
\]
This extracts Newton's second law from the system's assumptions; it is
the $P_6$ clause of the axiomatic presentation
\cite{mckinsey1953axiomatic}. In \texttt{FakeScience.lean}, the same
name is instead a theorem of type \texttt{False}. Lean rejects both a
joint import and the combined declarations because the name is already
declared. This diagnostic is a \emph{name collision}, not a theorem that
the two physical claims contradict each other.

The fake file also declares
\[
  \texttt{axiom ClassicalMechanics.physics : False}.
\]
This is a different and deeper problem: the file assumes a proof of
\texttt{False}. From that assumption Lean's \texttt{False.elim} derives
any proposition $Q$, so the resulting theory is logically inconsistent
relative to ordinary Lean logic. Lean accepts the file because it permits
users to introduce axioms; acceptance does not validate those axioms.
The fake theorem explicitly depends on \texttt{physics}; the genuine
\texttt{secondLawHolds} has no declared axiom dependencies. The supplied
fake file does not state an alternative force law or supply a model, so
no specific physical conflict with $P_6$ can be inferred from it.
Renaming the duplicate theorem would resolve the name collision, but
would not repair the false axiom.

\section{Reproduction}

With Lean 4 installed, run from the repository root:
\begin{verbatim}
  bash _lean/check_compatibility.sh
\end{verbatim}
The script compiles both files separately, checks the axiom dependencies,
and verifies that joint import and combined declarations fail with the
expected name-collision diagnostics. It also checks
\texttt{FakeScienceAudit.lean}, which derives an arbitrary proposition
from the fake axiom and reports the dependency using \texttt{\#print axioms}.
The checks reported here were run with Lean 4.35.0; the script checks
specific diagnostic wording, which may change between Lean versions.
On a machine using \texttt{elan} without a default toolchain, set
\texttt{ELAN\_TOOLCHAIN} to an installed Lean 4 toolchain first.

These checks do not prove that the abstract classical mechanics
signature has a concrete real-analysis interpretation, nor do they check
the truth of any empirical paper. They locate the collision and show
exactly how the fake file's explicit assumption makes arbitrary
propositions derivable.

\bibliographystyle{plain}
\bibliography{references}
\end{document}
